% Standalone XeLaTeX source; source problem: 2012ap_exam.pdf, PDF p. 7.
\documentclass[11pt,a4paper]{article}
\usepackage[margin=24mm]{geometry}
\usepackage{fontspec}
\usepackage{xeCJK}
\setCJKmainfont{HaranoAjiMincho-Regular.otf}
\setCJKsansfont{HaranoAjiGothic-Medium.otf}
\usepackage{amsmath,amssymb}
\setlength{\parindent}{1em}
\setlength{\parskip}{0.35em}
\title{物理工学専攻 院試解答集 ver2\\2012年 物理学II 第1問}
\author{政岡凜太郎}
\date{}
\begin{document}
\maketitle

生成・消滅演算子を用いると
\begin{equation}
  \hat X=x_0(\hat a+\hat a^\dagger),\qquad
  x_0=\sqrt{\frac{\hbar}{2m\omega_0}},\qquad
  \hat H_0=\hbar\omega_0\left(\hat a^\dagger\hat a+\frac12\right).
\end{equation}

\section*{[1] 双極子遷移}
$\hat\mu=q\hat X$ だから、数状態間の行列要素は
\begin{equation}
  \langle l|\hat\mu|n\rangle
  =q x_0\left(\sqrt n\,\delta_{l,n-1}
        +\sqrt{n+1}\,\delta_{l,n+1}\right).
\end{equation}
従って $n\ge1$ なら $|n-1\rangle,|n+1\rangle$ へ遷移でき、基底状態 $|0\rangle$ からは $|1\rangle$ だけが許される。

\section*{[2] 四次摂動によるエネルギー補正}
一次摂動論では $\Delta E_n=\langle n|\lambda\hat X^4|n\rangle$ である。
$Q=\hat a+\hat a^\dagger$ とおき、問題文の昇降作用を2回用いると
\begin{equation}
  Q^2|n\rangle=\sqrt{n(n-1)}|n-2\rangle
  +(2n+1)|n\rangle+\sqrt{(n+1)(n+2)}|n+2\rangle.
\end{equation}
$n<2$ のとき最初の項はなく、負の添字の状態は含めない。
$Q$ はHermiteであり、異なる数状態は直交するので、
\begin{align}
  \langle n|Q^4|n\rangle
  &=\|Q^2|n\rangle\|^2\\
  &=n(n-1)+(2n+1)^2+(n+1)(n+2)\\
  &=6n^2+6n+3.
\end{align}
したがって
\begin{equation}
  \boxed{\Delta E_n
  =3\lambda(2n^2+2n+1)\left(\frac{\hbar}{2m\omega_0}\right)^2.}
\end{equation}

\section*{[3] 摂動後の双極子遷移}
摂動を一次まで含む規格化固有状態は
\begin{equation}
 |n\rangle'=|n\rangle+
 \frac{\lambda x_0^4}{\hbar\omega_0}
 \sum_{\substack{k\ne n\\ |k-n|\in\{2,4\}}}
 \frac{\langle k|Q^4|n\rangle}{n-k}|k\rangle+O(\lambda^2),
 \qquad Q=\hat a+\hat a^\dagger.
\end{equation}
四次演算子は偶数準位だけを混ぜるので、$|8\rangle'$ と反対パリティの状態だけが双極子遷移の候補となる。さらに $Q$ は準位を1つだけ変えるため、候補は $l=3,5,7,9,11,13$ である。ただし候補範囲だけでは行列要素の非零性は決まらない。$Q_{lk}=\langle l|Q|k\rangle$ と書く。$Q$ を4回作用させると、非対角行列要素は
\begin{align}
 (Q^4)_{n+2,n}&=(4n+6)\sqrt{(n+1)(n+2)},\\
 (Q^4)_{n+4,n}&=\sqrt{(n+1)(n+2)(n+3)(n+4)}
 \qquad(n\ge0)
\end{align}
となる。準位を下げる行列要素は、これらの実行列要素の添字を交換して得られる。$\eta=\lambda x_0^4/(\hbar\omega_0)$ とすると、一次遷移振幅は
\begin{align}
 \frac{\langle l'|q\hat X|8'\rangle}{q x_0}
 &=Q_{l8}+\eta\left[
 \sum_{k\ne l}\frac{(Q^4)_{kl}Q_{k8}}{l-k}
 +\sum_{k\ne8}\frac{Q_{lk}(Q^4)_{k8}}{8-k}\right]+O(\eta^2).
\end{align}
$l=13$ では一次補正に寄与するのは、終状態の $|9\rangle$ 成分と初期状態の $|12\rangle$ 成分だけである。その二項は
\begin{equation}
  \frac{\sqrt{13}\,(Q^4)_{12,8}}{4}
  +\frac{\sqrt{13}\,(Q^4)_{12,8}}{-4}=0.
\end{equation}
上の角括弧内の二つの総和へ行列要素を代入する。各行で左の加数が終状態の補正、右の加数が初期状態の補正からの寄与である。
\begin{align}
 l=3:&\quad -2\sqrt{105}+2\sqrt{105}=0,\\
 l=5:&\quad -61\sqrt{21}+65\sqrt{21}=4\sqrt{21},\\
 l=7:&\quad -153\sqrt8+105\sqrt8=-48\sqrt8,\\
 l=9:&\quad 408-570=-162,\\
 l=11:&\quad 69\sqrt{110}-66\sqrt{110}=3\sqrt{110}.
\end{align}
従って $l=3$ でも一次振幅が消える。残る4つの振幅は
\begin{align}
 \frac{\langle5'|q\hat X|8'\rangle}{q x_0}&=4\sqrt{21}\,\eta+O(\eta^2),&
 \frac{\langle7'|q\hat X|8'\rangle}{q x_0}&=\sqrt8(1-48\eta)+O(\eta^2),\\
 \frac{\langle9'|q\hat X|8'\rangle}{q x_0}&=3(1-54\eta)+O(\eta^2),&
 \frac{\langle11'|q\hat X|8'\rangle}{q x_0}&=3\sqrt{110}\,\eta+O(\eta^2).
\end{align}
従って求める遷移先は
\begin{equation}
  \boxed{|5\rangle',\ |7\rangle',\ |9\rangle',\ |11\rangle'.}
\end{equation}

\section*{[4] Heisenberg表示の位置演算子}
$[\hat H_0,\hat a]=-\hbar\omega_0\hat a$、$[\hat H_0,\hat a^\dagger]=\hbar\omega_0\hat a^\dagger$ より
\begin{equation}
  \hat a(t)=e^{-i\omega_0t}\hat a(0),\qquad
  \hat a^\dagger(t)=e^{i\omega_0t}\hat a^\dagger(0).
\end{equation}
従って
\begin{align}
  \hat X(t)
  &=x_0\left(e^{-i\omega_0t}\hat a(0)
           +e^{i\omega_0t}\hat a^\dagger(0)\right)\\
  &=\boxed{\hat X(0)\cos(\omega_0t)
      +\frac{\hat P(0)}{m\omega_0}\sin(\omega_0t).}
\end{align}

\section*{[5] 位置の期待値}
数状態では $\langle n|\hat X(0)|n\rangle=\langle n|\hat P(0)|n\rangle=0$ なので
\begin{equation}
  \boxed{\langle n|\hat X(t)|n\rangle=0.}
\end{equation}
有限振幅で振動する例は、$\hat a$ の固有状態である規格化コヒーレント状態
\begin{equation}
  |\alpha\rangle=e^{-|\alpha|^2/2}
  \sum_{n=0}^{\infty}\frac{\alpha^n}{\sqrt{n!}}|n\rangle,
  \qquad \hat a|\alpha\rangle=\alpha|\alpha\rangle
\end{equation}
である。この状態では
\begin{equation}
  \boxed{\langle\alpha|\hat X(t)|\alpha\rangle
  =2x_0\operatorname{Re}\!\left(\alpha e^{-i\omega_0t}\right),}
\end{equation}
となり、$\alpha\ne0$ なら有限の振幅で角振動数 $\omega_0$ の正弦振動をする。

\end{document}
